% GENERATED FILE. Edit canonical lessons, then run npm run generate:books.
% canonical-source-hash: fnv1a64:47b1a1ed4ecce195
% canonical-lessons: FA-W23-he-jimi, FA-W23-jim, FA-W23-khe, FA-W23-che

\chapter{The Hook Family}
\label{ch:fa-letters-hook-family}
\hlchaptermodality{23}

\begin{chapteropening}
\textbf{By the end of this chapter:} \emph{I can write he-ye jimi, jim, khe and che, and tell the four apart by their dots.}

Complete the last lesson of chapter 23: i can write he-ye jimi, jim, khe and che, and tell the four apart by their dots.
\end{chapteropening}

\section[he-ye jimi]{\fa{ح} (he-ye jimi) — the letter that opens \fa{حال}}

\label{lesson:FA-W23-he-jimi}



\begin{quote}
Before the new one: write \fa{ز} and \fa{د} once each, and name them.

Now say \textbf{\fa{حال}} (\emph{hâl}), \textquotedblleft{}state, how you are\textquotedblright{}. You have read it for a long time. This lesson takes one letter out of it and puts it in your hand.
\end{quote}

\subsection*{Script you'll notice: \fa{ح}}
\textbf{\fa{ح}} — \emph{he-ye jimi}, the sound \emph{h}.

It opens \textbf{\fa{حال}} (\emph{hâl}), the \textquotedblleft{}how you are\textquotedblright{} in \textbf{\fa{حالت} \fa{چطور} \fa{است؟}}. Persian already has an \emph{h}, \textbf{\fa{ه}}, which you write; this second one came in with Arabic words and sounds the same. Its name, \emph{he-ye jimi}, means \textquotedblleft{}the \emph{h} shaped like \emph{jim}\textquotedblright{}: a short head and a deep bowl, with \textbf{no dots}.

\subsection*{Writing: \fa{ح}}
\hlblockfigure{\detokenize{figures/FA-W23-he-jimi-filmstrip.pdf}}{How \fa{ح} is written, stroke by stroke}

\begin{itemize}
\raggedright
  \item \textbf{1.} start here: draw the short upper head from left to right
  \item \textbf{2.} continue down and around the deep bowl without lifting — and only now lift
\end{itemize}

\textbf{Pen lifts: 0.} The pen never leaves the paper.

\begin{quote}
Stroke order is one attested teaching order; hands and teachers vary. Source: Persian Online, How to Write Persian Characters, \fa{ح} demonstration at 00:42–00:46 (Liberal Arts Instructional Technology Services, Univ. of Texas at Austin).
\end{quote}

\subsection*{Your turn}
\begin{itemize}
\raggedright
  \item \emph{Look:} at \textbf{\fa{حال}} and point at \fa{ح} inside it
  \item \emph{Trace it:} \fa{ح} three times, saying \emph{h} as you finish each one
  \item \emph{Write it:} \fa{د}, then \fa{ح}, from memory
  \item \emph{Read it:} \textbf{\fa{حال}} aloud once more
\end{itemize}

\begin{tcolorbox}[breakable,colback=teal!4,colframe=teal!35!black,title={Before you move on}]
Which letter is this — \fa{ح}? (\textbf{he-ye jimi}, \emph{h}.) Which word did it come from? (\textbf{\fa{حال}}, \emph{hâl}, \textquotedblleft{}state, how you are\textquotedblright{}.)
\end{tcolorbox}
\section[jim]{\fa{ج} (jim) — the letter that ends \fa{پنج}}

\label{lesson:FA-W23-jim}



\begin{quote}
Before the new one: write \fa{د} and \fa{ح} once each, and name them.

Now say \textbf{\fa{پنج}} (\emph{panj}), \textquotedblleft{}five\textquotedblright{}. You have read it for a long time. This lesson takes one letter out of it and puts it in your hand.
\end{quote}

\subsection*{Script you'll notice: \fa{ج}}
\textbf{\fa{ج}} — \emph{jim}, the sound \emph{j}.

It ends \textbf{\fa{پنج}} (\emph{panj}), \textquotedblleft{}five\textquotedblright{}. It is \textbf{\fa{ح}}, from the last lesson, with \textbf{one dot below} the head.

\subsection*{Writing: \fa{ج}}
Put your pen on \fa{ج} and follow its line. Copy the shape you can see — slowly, and larger than it is printed.

\begin{quote}
This book does not yet tell you \textbf{where to start this letter or which way to travel}. It is not written down here until it can be written down with a source. Copying what is in front of you needs no such source, and it is how the shape gets into your hand in the meantime.
\end{quote}

\subsection*{Your turn}
\begin{itemize}
\raggedright
  \item \emph{Look:} at \textbf{\fa{پنج}} and point at \fa{ج} inside it
  \item \emph{Trace it:} \fa{ج} three times, saying \emph{j} as you finish each one
  \item \emph{Write it:} \fa{ح}, then \fa{ج}, from memory
  \item \emph{Read it:} \textbf{\fa{پنج}} aloud once more
\end{itemize}

\begin{tcolorbox}[breakable,colback=teal!4,colframe=teal!35!black,title={Before you move on}]
Which letter is this — \fa{ج}? (\textbf{jim}, \emph{j}.) Which word did it come from? (\textbf{\fa{پنج}}, \emph{panj}, \textquotedblleft{}five\textquotedblright{}.)
\end{tcolorbox}
\section[khe]{\fa{خ} (khe) — the letter that opens \fa{خدا}}

\label{lesson:FA-W23-khe}



\begin{quote}
Before the new one: write \fa{ح} and \fa{ج} once each, and name them.

Now say \textbf{\fa{خدا}} (\emph{khodâ}), \textquotedblleft{}God\textquotedblright{}. You have read it for a long time. This lesson takes one letter out of it and puts it in your hand.
\end{quote}

\subsection*{Script you'll notice: \fa{خ}}
\textbf{\fa{خ}} — \emph{khe}, a \emph{kh} scraped at the back of the throat.

It opens \textbf{\fa{خدا}} (\emph{khodâ}). The same head and bowl as \textbf{\fa{ح}}, with \textbf{one dot above}. One skeleton, and the dot decides the sound.

\subsection*{Writing: \fa{خ}}
\hlblockfigure{\detokenize{figures/FA-W23-khe-filmstrip.pdf}}{How \fa{خ} is written, stroke by stroke}

\begin{itemize}
\raggedright
  \item \textbf{1.} start here: draw the short upper head from left to right
  \item \textbf{2.} without lifting, continue down and around the deep bowl
  \item \textbf{3.} lift once, then place the dot above — and only now lift
\end{itemize}

\textbf{Pen lifts: 1.}

\begin{quote}
Stroke order is one attested teaching order; hands and teachers vary. Source: Persian Online, How to Write Persian Characters, \fa{خ} demonstration at 00:49–00:54 (Liberal Arts Instructional Technology Services, Univ. of Texas at Austin).
\end{quote}

\subsection*{Your turn}
\begin{itemize}
\raggedright
  \item \emph{Look:} at \textbf{\fa{خدا}} and point at \fa{خ} inside it
  \item \emph{Trace it:} \fa{خ} three times, saying \emph{kh} as you finish each one
  \item \emph{Write it:} \fa{ج}, then \fa{خ}, from memory
  \item \emph{Read it:} \textbf{\fa{خدا}} aloud once more
\end{itemize}

\begin{tcolorbox}[breakable,colback=teal!4,colframe=teal!35!black,title={Before you move on}]
Which letter is this — \fa{خ}? (\textbf{khe}, \emph{kh}.) Which word did it come from? (\textbf{\fa{خدا}}, \emph{khodâ}, \textquotedblleft{}God\textquotedblright{}.)
\end{tcolorbox}
\section[che]{\fa{چ} (che) — the letter that opens \fa{چیست}}

\label{lesson:FA-W23-che}



\begin{quote}
Before the new one: write \fa{ج} and \fa{خ} once each, and name them.

Now say \textbf{\fa{چیست}} (\emph{chist}), \textquotedblleft{}what is?\textquotedblright{}. You have read it for a long time. This lesson takes one letter out of it and puts it in your hand.
\end{quote}

\subsection*{Script you'll notice: \fa{چ}}
\textbf{\fa{چ}} — \emph{che}, the sound \emph{ch}.

It opens \textbf{\fa{چیست}} (\emph{chist}), \textquotedblleft{}what is?\textquotedblright{}, one of your first questions. It is one of the four letters Persian \textbf{added} to the Arabic alphabet: the same skeleton as \textbf{\fa{ح}}, with \textbf{three dots below}.

\subsection*{Writing: \fa{چ}}
\hlblockfigure{\detokenize{figures/FA-W23-che-filmstrip.pdf}}{How \fa{چ} is written, stroke by stroke}

\begin{itemize}
\raggedright
  \item \textbf{1.} start here: draw the short upper head from left to right
  \item \textbf{2.} continue down and around the deep bowl without lifting
  \item \textbf{3.} lift, then place the lower-left dot
  \item \textbf{4.} lift again and place the lower-right dot
  \item \textbf{5.} lift again and place the lower-center dot — and only now lift
\end{itemize}

\textbf{Pen lifts: 3.}

\begin{quote}
Stroke order is one attested teaching order; hands and teachers vary. Source: Persian Online, How to Write Persian Characters, \fa{چ} demonstration at 00:35–00:41 (Liberal Arts Instructional Technology Services, Univ. of Texas at Austin).
\end{quote}

\subsection*{Your turn}
\begin{itemize}
\raggedright
  \item \emph{Look:} at \textbf{\fa{چیست}} and point at \fa{چ} inside it
  \item \emph{Trace it:} \fa{چ} three times, saying \emph{ch} as you finish each one
  \item \emph{Write it:} \fa{خ}, then \fa{چ}, from memory
  \item \emph{Read it:} \textbf{\fa{چیست}} aloud once more
\end{itemize}

\begin{tcolorbox}[breakable,colback=teal!4,colframe=teal!35!black,title={Before you move on}]
Which letter is this — \fa{چ}? (\textbf{che}, \emph{ch}.) Which word did it come from? (\textbf{\fa{چیست}}, \emph{chist}, \textquotedblleft{}what is?\textquotedblright{}.)
\end{tcolorbox}
